\section{Test Case 13: dynamic-topology GPU baseline}

Test Case 13 initializes 1,024 beads uniformly between the nozzle and
collector and enables both \texttt{inserting yes} and \texttt{removing yes}.
It retains the material parameters and constant axial field of Test Case 9.
An intentionally artificial nozzle velocity of 2,000 cm/s forces multiple
topology events during a short 1,000-step run. Its input is stored in
\texttt{examples/input-13/input.dat}.

Both call-scoped and persistent executions contain 13 insertions and 13
removals, finish with 1,024 active beads, and remain between 1,023 and 1,025
active beads. A
structured \texttt{Topology event:} line is printed whenever insertion or
removal occurs.

The current dynamic milestone reserves capacity for 1,280 beads and keeps the
RK4 state, Coulomb data, derivatives, and scratch arrays persistently mapped
while the active lower and upper bead indices change. Collector detection and
clamping run on the device and return one decision scalar. Nozzle distance
checks, blocked-bead release, new-record initialization, and the
\texttt{npjet} update also execute on the device. The host receives topology
scalars each step and the two new tail records only when insertion occurs.
Removed records are synchronized
individually for output and cleared on both host and device. Reallocation
beyond the reserved capacity and general compaction remain future work. An
initial call-scoped NVFORTRAN 24.3 and NVIDIA A30 measurement required
40.659696 s on CPU and 3.834931 s with OpenACC. The bounded persistent path
with host insertion required 2.829376 s. Device-side insertion reduces this
to 2.623558 s, or 381.162 steps/s. Moving RK4 intermediate updates from
host to device changes floating-point rounding. The insertion threshold
amplifies this into progressively earlier additions, while addition and
removal counts, active bounds, and final bead count remain unchanged.

With the environment variable
\texttt{JETSPIN\_TOPOLOGY\_SNAPSHOT=1}, JETSPIN writes the complete active
state at every topology event. This diagnostic showed identical bead indices,
frozen flags, masses, charges, volumes, and event metadata. The first device
difference appears in a transverse quantity close to zero at step 40. Running
the same OpenACC equation-of-motion kernel on the CPU agreed with the original
CPU implementation near machine precision. The observed growth therefore
originates primarily from device floating-point evaluation of curvature and
is not evidence of an insertion or removal indexing error. The persistent
dynamic control path retains all 26 events but has a separate versioned A30
event stream. Setting
\texttt{JETSPIN\_OPENACC\_DISABLE\_PERSISTENT=1} restores the call-scoped
path and exactly restores the previous event timesteps on the same A30. A
separate diagnostic using the historical host cross-section computation did
not change the persistent stream, isolating the timing shift to GPU RK4
arithmetic rather than insertion, removal, or Coulomb indexing.
Moving the insertion distance test to the GPU shifts one later threshold
crossing from step 910 to step 909, without changing event counts, active
bounds, or final topology.
